\documentclass[12pt,draftcls]{ucdavisthesis}

% PLEASE READ THE MANUAL - ucdavisthesis.pdf (in the package installation directory)

%%%%%%%%%%%%%%%%%%%%%%%%%%%%%%%%%%%%%%%%%%%%%%%%%%%%%%%%%%%%%%%%%%%%%%%%
%                                                                      %
%               LATEX COMMANDS FOR DOCUMENT SETUP                      %
%                                                                      %
%%%%%%%%%%%%%%%%%%%%%%%%%%%%%%%%%%%%%%%%%%%%%%%%%%%%%%%%%%%%%%%%%%%%%%%%

%\usepackage{bookmark}
\usepackage[us,nodayofweek,12hr]{datetime}
\usepackage{graphicx}
%\usepackage[square,comma,numbers,sort&compress]{natbib}
%\usepackage{hypernat}
% Other useful packages to try
%\usepackage{amsmath}
%\usepackage{amssymb}
%
% Different fonts to try (uncomment only fontenc and one font at a time)
% (you may need to install these first)
%\usepackage[T1]{fontenc} %enable fontenc package if using one of the fonts below
%\usepackage[adobe-utopia]{mathdesign}
%\usepackage{tgschola}
%\usepackage{tgbonum}
%\usepackage{tgpagella}
%\usepackage{tgtermes}
%\usepackage{fourier}
%\usepackage{fouriernc}
%\usepackage{kmath,kerkis}
%\usepackage{kpfonts}
%\usepackage[urw-garamond]{mathdesign}
%\usepackage[bitstream-charter]{mathdesign}
%\usepackage[sc]{mathpazo}
%\usepackage{mathptmx}
%\usepackage[varg]{txfonts}

\hyphenation{dis-ser-ta-tion blue-print man-u-script pre-par-ing} %add hyphenation rules for words TeX doesn't know


%\renewcommand{\rightmark}{\scriptsize A University of California Davis\ldots \hfill Rev.~\#1.0 \quad Compiled: \currenttime, \today}
% a fancier running header that can be used with draftcls options

%%%%%%%%%%%%%%%%%%%%%%%%%%%%%%%%%%%%%%%%%%%%%%%%%%%%%%%%%%%%%%%%%%%%%%%%
%                                                                      %
%        DOCUMENT SETUP AND INFORMATION FOR PRELIMINARY PAGES          %
%                                                                      %
%%%%%%%%%%%%%%%%%%%%%%%%%%%%%%%%%%%%%%%%%%%%%%%%%%%%%%%%%%%%%%%%%%%%%%%%

\title          {A University of California Davis\\
                 Dissertation/Thesis LaTeX Class File}
%Exact title of your thesis. Indicate italics where necessary by underlining or using italics. Please capitalize the first letter of each word that would normally be capitalized in a title.

\author         {Dustin Ray Burns}
%Your full name as it appears on University records. Do not use initials.

\authordegrees  {B.S. (First Institution) 2002 \\
                 M.S. (Second Institution) 2004}
%Indicate your previous degrees conferred.

\officialmajor  {Official Name of Major}
%This is your official major as it appears on your University records.

\graduateprogram{Graduate Program Name}
%This is your official graduate program name. Used for UMI abstract.

\degreeyear     {2020}
% Indicate the year in which your degree will be officially conferred.

\degreemonth    {June}
% Indicate the month in which your degree will be officially conferred. Used for UMI abstract.

\committee{Committee Chair Name, Chair}{Committee Member Name}{Committee Member Name}{}{}
% These are your committee members. The command accepts up to five committee members so be sure to have five sets of braces, even if there are empties.

%%%%%%%%%%%%%%%%%%%%%%%%%%%%%%%%%%%%%%%%%%%%%%%%%%%%%%%%%%%%%%%%%%%%%%%%

%\copyrightyear{2020}
%\nocopyright

%%%%%%%%%%%%%%%%%%%%%%%%%%%%%%%%%%%%%%%%%%%%%%%%%%%%%%%%%%%%%%%%%%%%%%%%

\dedication{\textsl{To someone very important \ldots \\
            a nice dedication.}}

%%%%%%%%%%%%%%%%%%%%%%%%%%%%%%%%%%%%%%%%%%%%%%%%%%%%%%%%%%%%%%%%%%%%%%%%

\abstract{The abstract submitted as part of your dissertation, in the introductory pages, does not have a word limit. It should follow the same format as the rest of your dissertation (1.5 inch left margin, double-spaced, consecutive page numbering, etc.).}

%%%%%%%%%%%%%%%%%%%%%%%%%%%%%%%%%%%%%%%%%%%%%%%%%%%%%%%%%%%%%%%%%%%%%%%%

\acknowledgments{Acknowledgements to those who helped you get to this point. They should be listed by chapter when appropriate.}

%%%%%%%%%%%%%%%%%%%%%%%%%%%%%%%%%%%%%%%%%%%%%%%%%%%%%%%%%%%%%%%%%%%%%%%%

% Each chapter can be in its own file for easier editing and brought in with the \include command.
% Then use the \includeonly command to speed compilation when working on a particular chapter.
%%% \includeonly{ucdavisthesis_example_Chap1}

\begin{document}

\newcommand{\bibfont}{\singlespacing}
% need this command to keep single spacing in the bibliography when using natbib

\bibliographystyle{unsrtnat}
%many other bibliography styles are available (IEEEtran, mla, etc.). Use one appropriate for your field.

\makeintropages %Processes/produces the preliminary pages

\chapter{Information on Preparation of a Dissertation or Thesis}
%
The text for this example dissertation is taken from the Graduate Studies web site\footnote{Please see \texttt{http://gradstudies.ucdavis.edu/students/filing.html} for the latest information.}. Please read through it carefully to make sure that your dissertation/thesis meets the university requirements.

For information and assistance in the preparation of your thesis or dissertation, please contact the Student Affairs Coordinator responsible for your program. You must make an appointment with the Student Affairs Coordinator in order to file your thesis or dissertation.

You are responsible for observing the filing dates and for preparing the thesis in the proper format. The thesis or dissertation must be accompanied by either the University Library Questionnaire (master's students) or the UMI Doctoral Dissertation Agreement form (Ph.D. students) on which you indicate your willingness either to have the University supply copies of your thesis to interested persons immediately, or that such permission should be withheld for a period of time, up to as maximum of three years.

Filing your dissertation or thesis is the last step in the process leading to the awarding of your degree. The final copy of your thesis or dissertation, which is ultimately deposited in the University Library, becomes a permanent and official record. If your manuscript is a doctoral dissertation, you have the choice of submitting either your entire manuscript or only your dissertation abstract to University Microfilms International (UMI) Dissertation Publishing. Both options provide higher visibility of your achievements. If you submit your entire dissertation to UMI, it will be returned to Shields Library where it will be bound and stored. The packet of information that was sent to you when we returned your approved candidacy form contained the forms you need to submit with your manuscript.


\section{Specifications}
%
Consult the calendar for the dates to file your thesis or dissertation at Graduate Studies. The dates are also printed in the General Catalog. You are responsible for observing the filing dates and for preparing the thesis or dissertation in the proper form.

\subsection{Style and Format}
%
You should be guided on matters of style by the chair and members of your thesis committee. Graduate Studies is not concerned with the form of the bibliography, appendix, footnotes, etc. as long as they are done in some acceptable, consistent and recognized manner approved by your committee.

There are many valuable references available to assist students in preparing and writing research papers and theses \cite{Gibaldi:80, Strunk:00, UoCPress:03}. Listed below are references that have been suggested by students and faculty\footnote{Be sure to search the internet for additional writing resources.}.
%
\begin{itemize}
\item MLA Handbook for Writers of Research Papers, Theses and Dissertations, by Joseph Gibaldi
    and Walter S. Achtert: The Modern Language Association of America, 4th ed., 1980.
\item The Elements of Style, by William Strunk, Jr., E.B. White and Roger Angell: Longman, 4th ed., 2000.
\item The Lively Art of Writing, by Lucile V. Payne, Mentor Books, reissue ed., 1983.
\item The Chicago Manual of Style, University of Chicago Press, 15th ed., 2003.
\end{itemize}

\subsection{Margins}
%
Every page of the dissertation or thesis must be kept within these margins:

\noindent left - 1.5 inches (for binding purposes and subsequent copying)\\
top, bottom, right - one inch\\
(page number may be outside of these margins)

Please Note: These minimum specifications also apply to all figures, charts, graphs, illustrations and appendices. Any pages submitted with less than the minimum margins will be returned to you.

\subsection{Spacing}
%
Double spacing should be used in typing the thesis, except in those places where conventional usage calls for single spacing -- for example, footnotes, indented quotations, tables, and the bibliography.

\subsection{Pagination}
%
Do not start renumbering pages at any point in your thesis or dissertation.

The preliminary pages, including the title page are numbered with small Roman Numerals which are centered at the bottom of the page. Begin numbering the preliminaries in lower case roman numerals with the title page which is always. The text and all other pages of the thesis or dissertation, including charts, figures, caption pages, maps and appendices are all numbered consecutively starting with Arabic 1 in the upper right hand corner about 1/2 inch from each edge. (In most cases this will start with the Introduction or Chapter 1.) If material is inserted after printing, it should be numbered 21a, 21b, and so on if it follows page 21. If a numbered page is later removed, a numbered page, blank except for the note saying that the original page had been deleted, should be inserted in the proper place. Any blank pages you insert as separators need to be numbered and marked as ``BLANK'' so that there is no confusion over missing information.

EVERY PAGE MUST BE NUMBERED CONSECUTIVELY.

\subsection{Typeface}
%
Consistency is essential. Any typeface style is acceptable, but the same typeface size and style must be used throughout. Use a font size between 10 and 13 points. Be sure your computer printer output is clean and dark. Any symbols, equations, etc.\ that are hand-drawn must be done in permanent black ink.

\subsection{Paper}
%
The paper selected must meet the following criteria: white in color, of standard size $8.5\,\times\,11$~inches, unpunched and of good quality. YOU ARE REQUIRED to use at least 25\% cotton bond paper that is at least 20~lb.\ weight; please select paper with a watermark so that it is apparent that the paper meets the minimum requirements. (Acid-free paper is suggested, but not required.) Only one side of the paper may be used for printing. The use of erasable bond paper or any copier type paper is not permitted. Any thesis or dissertation submitted on unacceptable paper will be returned to the student.

Exceptions regarding size and type of paper may be made in certain fields, for example, music, art, and geology. If you are planning to submit a thesis\slash dissertation that does not meet the standard requirements, you should consult Graduate Studies before preparing in final form. 
\chapter{The Second Chapter}

\section{Overall Appearance}
%
Do not bind or punch holes in the thesis or dissertation; the document will be accepted as loose pages only.

The general appearance of the manuscript is important, not only for Library use at UC Davis, but also since the copy which doctoral students submit will be sent to UMI Dissertation Publishing to be copied on request for scholars around the world. For this reason, correction of typographical errors should be made with care. Pages with interlineations, crossed out words or letters, strike-overs and unsightly corrections will be rejected and the candidate will be required to replace them. The use of white correction fluid or white correction carbon IS NOT ACCEPTABLE. If you photocopy your dissertation or thesis from laser/copier paper to the accepted 25\% cotton bond thesis paper, please make certain each page is clear, with no dark edges.

\subsection{Tables, Graphs and Captions}
%
Charts and tables may be placed horizontally or vertically, but in either case must fit within the required margins. It may be necessary to use a reducing copier in order to achieve this. If necessary, wide tables, charts, and figures can be placed sideways. Figures may be embodied in the text or take up a full page. Each figure or table must be numbered consecutively (do not renumber each chapter unless you include chapter numbers, e.g., Fig.~1.1, Fig.~2.1, etc.) and should have a caption. With full page figures, this can be done on the facing caption page. Such a caption page must be the mirror image of a normal page, i.e., the wider margin is on the right and the number in the upper left hand corner. The caption itself should be centered and should always be on a page by itself, not on the back of a preceding page.

NOTE: If your figures or charts are placed horizontally on the page, your page number must still be in the upper right corner when the page is viewed in the manuscript. Pagination must be consistent.

\subsection{Photographs, Illustrations and Maps}
%
Plates or figures should be the original photographs or drawings, or photographic reproductions of the originals. All reproductions must be clean and distinct. Photographs should be positive prints; line drawings or tabular material reproduced by photography should be positives. If color is important, special markings should be used to identify color variations. If a plate is smaller than the standard $8.5\times 11$~inch page, it should be attached to a blank page of standard size with dry mounting tissue, white glue or a permanent spray adhesive. If white glue is used, the complete perimeter should be covered. Do not use rubber cement, tapes, mucilage, or photo-corners to attach plates. These are cited because they are well known and are not permanent processes; this is not a complete list of unacceptable methods. You should inquire about the acceptability of any other method that you may consider using. If a digital graphic is used and printed on photographic paper, this is acceptable if used only in this manner. Pagination must be consistent.

\subsection{Oversized Material}
%
Tables, charts, and plates that are larger than $8.5\times 11$~inches in size must not be folded, since they will have to be folded in a special manner for binding. They should be submitted in a cardboard cylinder, available at the UC Davis Bookstore or any stationery store, and must also be numbered with Arabic numbers.

Extremely oversized material should also be submitted in a cardboard cylinder. These oversized sheets should not be assigned page numbers, but your name and title of your dissertation or thesis should appear on each sheet in the lower right hand corner. The Library will fold this material for placement in a special pocket when the manuscript is bound.

\subsection{Using Published Material}
%
If approved by the thesis or dissertation committee, reports of research undertaken during graduate study at UC Davis which have been published may be accepted in printed form as all or part of the master's thesis or doctoral dissertation. If you are not the sole or first author of the published material submitted, the use of co-authored materials must be approved by the department or graduate group concerned.

The pages of the published material must have the same margins and type of paper as specified in these instructions which may necessitate reduction of the material. Note that enlargement of the materials to meet standard margin requirements is not needed. When using reprints, page numbers of the reprint should be removed and replaced with numbers corresponding to the position of the reprint within the dissertation. Each chapter may have an abstract of its own. For doctoral students there must be a general abstract covering the entire dissertation.

If the papers or articles to be submitted are not yet in published form, they must conform to the specifications of these instructions.

\subsection{Computer-produced Dissertations and Theses}
%
Manuscripts done on a computer or word processor should be produced either on a hard-copy terminal or phototypesetter. These are just a few of the software packages which will operate within the format described in these instructions:

\begin{itemize}
  \item \TeX and \LaTeX
  \item Deltagraph
  \item Microsoft Word
  \item Wordperfect
\end{itemize}

Laser and other high quality printers are acceptable.


\subsection{Methods for Printing Final Copies}
%
Graduate Studies will accept a thesis or dissertation prepared by any one of a number of processes provided that it meets all of the following conditions:

\begin{itemize}
  \item it must be a legible, clean, clear copy;
  \item it must be photographically reproducible;
  \item the ink must be permanent;
  \item a minimum of 20~lb.\ 25\% cotton bond paper is required. (Please select paper with a watermark so that it is apparent that the paper meets the minimum requirements.)
\end{itemize}

The following are acceptable processes: laser printer, ink jet printer, typed, Xerox, etc. This is not a complete list of acceptable processes and you should consult Graduate Studies concerning any process you are planning to use. Even though a process is acceptable on other grounds, the copy submitted will be rejected if it is not clear and legible. If you are uncertain about the acceptability of any copy you are planning to submit, please consult Graduate Studies.

The following processes are NOT acceptable: Ditto, Hectograph, Thermofax, blueprint, dot matrix printer. These are mentioned because they are widely used for other purposes; this is not a complete list of unacceptable processes.

Before submitting the copy to Graduate Studies, check to be sure that every page has been copied, numbered correctly, and is in the right order.

The ORIGINAL copy of the title page containing the original signatures of your committee must, in all cases, be submitted with the thesis or dissertation. It is preferable that the title page be printed on 25\% cotton bond paper.

\subsection{Number of Copies}
%
One unbound approved copy must be filed in Graduate Studies for later deposit in the University Library. You should consult your graduate staff person, graduate adviser, and committee chair to determine whether additional copies are required by your department or graduate group. You are responsible for providing any additional copies in these cases. Two places that offer binding are Cal Na Binding in Sacramento (916-447-4355) and Navin's Copy Shop in Davis (530-758-2311) which contracts with Cal Na.

\subsection{Dates of Filing}
%
Check the calendar for deadlines for filing the master's thesis or the doctoral dissertation with the committees in charge and with Graduate Studies. Deadlines are also announced each year in the Class Schedule and Registration Guide and the General Catalog. The deadline for filing with your committee is a recommended deadline to allow time for making revisions. The deadline for filing with Graduate Studies is firm.

It is important to bring all documents, forms and supplies with you when you file your thesis or dissertation. Please review the checklist for master's or doctoral students prior to your appointment.

\subsection{Title Page}
%
Graduate Studies does not supply the title page. You must prepare your title page in accordance with the sample. The title page is to be signed by all members of your committee when they have approved the thesis or dissertation. Only the original title page will be accepted with the thesis or dissertation.

\subsection{Dissertation Abstracts}
%
You must provide two different forms of your abstract. The first is submitted separately from the dissertation to University Microfilms International (UMI) Dissertation Publishing. The other format follows the regular formatting guidelines and is submitted as part of your dissertation.

The abstract that is submitted to UMI must be formatted as shown in the example here. The body of the abstract cannot exceed 350 words. It should be in typewritten form, double-spaced, and on bond paper. It is important to write an abstract that gives a clear description of the content and major divisions of the dissertation, since UMI will publish the abstract exactly as submitted. Students completing their requirements under Plan A should provide extra copies of the typed summary for use by the dissertation committee during the examination.

The abstract submitted as part of your dissertation, in the introductory pages, does not have a word limit. It should follow the same format as the rest of your dissertation (1.5 inch left margin, double-spaced, consecutive page numbering, etc.).

\subsection{University Microfilms International}
%
As a doctoral candidate, you will sign an agreement with University Microfilms International/Proquest Information and Learning Company (PQIL) for the microfilming of the dissertation or the printing of the abstract on the dissertation abstract database. Forms are available in the Office of Graduate Studies and should be filled out and submitted to Graduate Studies with the dissertation. You will be given the option to apply for copyright, but it is not required. If you would like to copyright your dissertation, please see the section on copyright and publication.

\subsection{Diploma}
%
When you file your thesis or dissertation, you will receive a temporary certificate that states you have completed all the requirements for your program and the official conferral date of your degree. This certificate may be given to your employer for proof of degree until the Registrar's Office issues an official transcript or diploma. You must complete a form to request your transcript or diploma. Official transcripts normally are available two months after the official degree conferral date; diplomas normally are available four months after this date.

\subsection{Copyright and Publication}
%
The copyright law of the United States is quite complex. The information contained in this section is only a general guide--more detailed information must be obtained from other sources. Whether or not you copyright your thesis or dissertation, you retain the right to publish all or any part of it by any means at any time, except for reproduction from a negative microfilm as described in the agreement form with University Microfilms International, if you signed such an agreement. Should you decide to copyright your thesis or dissertation, you must include a separate copyright page after the title page; do not use a page number on the copyright page. By adding this copyright notice, which should be included in all copies you distribute, you have copyrighted your thesis or dissertation. At this point you have several options:

If you are a doctoral student, you may have the copyright registered for you by roquest Information and Learning Company (PQIL). Along with the UMI Doctoral Dissertation Agreement form, you will need to submit a U.S. Postal money order in the amount of \$65 (no personal checks), payable to PQIL, to cover the copyright fee and the cost of two positive microfilm copies which will be filed with the Copyright Office.

You may register the copyright yourself by submitting to the Registrar of Copyrights the appropriate application form, a filing fee and one or two copies of the work. In order to have full protection against infringement, this should be done as soon as possible. Information and forms can be obtained from the Registrar of Copyrights, Library of Congress, Washington D.C. 20559.

You may choose to copyright your thesis or dissertation by adding the copyright notice, submitting a copy to the Registrar of Copyrights, but not registering it. (Federal copyright law requires that copies of all works published with notices of copyright be deposited with the Library of Congress, even if the copyright is not registered). However, to protect your rights in a copyright dispute and in order to be compensated for damages caused by infringement, your copyright must be registered.

\subsection{Permissions}
%
Since the depositing of your thesis or dissertation in the University Library and/or its being made available by UMI/PQIL may constitute a form of publication, you may have to obtain permission to use (or quote) copyrighted material, such as that in most journal articles or books. It is the author (i.e., you) who is responsible in the matter of copyrighted materials. The agreement form which you sign with University Microfilms specifically absolves them of any such responsibility.

If you quote extensively from a particular author, especially in fields such as fiction, drama, criticism, or poetry, or if copyrighted maps, charts, statistical tables, or similar materials have been reproduced, you must write the copyright owner(s), describe the use which you are making of the materials, and request permission to use it in the dissertation or thesis.

For your protection, a statement listing such materials should be included in the acknowledgements of the dissertation or thesis. The statement should inform the reader (1) that permission has been granted for their use, and (2) the source of the permission.

\subsection{Filing Fee}
%
The Filing Fee was established expressly to assist those students who have completed all requirements for degrees except filing theses or dissertations and/or taking final examinations (master's comprehensive exams or doctoral final examinations) and are no longer using University facilities. The Filing Fee is a reduced fee paid in lieu of registration fees. It is assessed only once and must be paid to the Cashier's Office prior to submission of the form to Graduate Studies. Filing Fee status restrictions (more restrictions are noted on the application instruction sheet):
%
\begin{itemize}
  \item You may not be using University facilities;
  \item You cannot be using faculty time other than the time involved in the final reading of the thesis or dissertation or in holding final examinations;
  \item You are not eligible to hold any academic appointment title for more than 1 quarter;
  \item You cannot hold a fellowship or receive financial aid.
\end{itemize}

If you are eligible to use the Filing Fee procedure, you should complete a Filing Fee application, obtain the signature of the Graduate Adviser and your Committee chair, and return the application to Graduate Studies before you stop registering. The Filing Fee must be paid prior to submitting the application to Graduate Studies.

Original (initial) filing fee deadlines adhere to registration deadlines. For example, if you were approved for one quarter of filing fee, you would be eligible to submit your dissertation or thesis up to the last day of late registration for the following quarter. If you do not submit by that date your filing fee status will expire and you would need to secure an extension from your program and from Graduate Studies. Filing fee extensions end the date noted on the petition. Make sure your filing fee is current before you submit your dissertation or thesis.

\chapter{Examples of Formatting and Commands}
%
The following chapter shows various formatting and sectioning commands. It also includes examples of how to insert figures and tables into the document. These examples are by no means comprehensive, they are just a sample of what can be done with the power of \LaTeX\@.

Filing your dissertation or thesis is the last step in the process leading to the awarding of your degree. The final copy of your thesis or dissertation, which is ultimately deposited in the University Library, becomes a permanent and official record.

\section[TOC Section Title]{First Section Title -- A Particularly Long Section Title with a Special Title for the TOC}
%
A short version of the section title can be included in square brackets just before the full title. This applies to subsections as well.

Filing your dissertation or thesis is the last step in the process leading to the awarding of your degree. The final copy of your thesis or dissertation, which is ultimately deposited in the University Library, becomes a permanent and official record.

\subsection[A Very Long Subsection Title to Test the TOC Wrapping Ability]{Subsection Title}
%
Filing your dissertation or thesis is the last step in the process leading to the awarding of your degree. The final copy of your thesis or dissertation, which is ultimately deposited in the University Library, becomes a permanent and official record.

Filing your dissertation or thesis is the last step in the process leading to the awarding of your degree. The final copy of your thesis or dissertation, which is ultimately deposited in the University Library, becomes a permanent and official record.

\subsubsection{Subsubsection Title}
%
Filing your dissertation or thesis is the last step in the process leading to the awarding of your degree. The final copy of your thesis or dissertation, which is ultimately deposited in the University Library, becomes a permanent and official record.

Filing your dissertation or thesis is the last step in the process leading to the awarding of your degree. The final copy of your thesis or dissertation, which is ultimately deposited in the University Library, becomes a permanent and official record.

\paragraph{Paragraph Title}
%
I have no idea who would use sectioning down to the paragraph level in a dissertation, but it is available for those who think it is necessary.

\section{Some Examples of Figures and Tables}
%
A table such as that shown in Table~\ref{tab.example_1} can be included easily. The hard part will be creating the table. Filing your dissertation or thesis is the last step in the process leading to the awarding of your degree. The final copy of your thesis or dissertation, which is ultimately deposited in the University Library, becomes a permanent and official record.
%
\begin{table}[tbh]
% increase table row spacing, adjust to taste
\renewcommand{\arraystretch}{1.3}
% if using array.sty, it might be a good idea to tweak the value of
% \extrarowheight as needed to properly center the text within the cells
\centering
% Some packages, such as MDW tools, offer better commands for making tables
% than the plain LaTeX2e tabular which is used here.
\begin{tabular}{|c||c|}
\hline
One & Two\\
\hline
Three & Four\\
\hline
\end{tabular}
\caption[Example table]{An example of a table. Notice the caption is centered except when it runs longer than a single line on the page.}
\label{tab.example_1}
\end{table}

Filing your dissertation or thesis is the last step in the process leading to the awarding of your degree. The final copy of your thesis or dissertation, which is ultimately deposited in the University Library, becomes a permanent and official record.
%
\begin{table}[tbh]
% increase table row spacing, adjust to taste
\renewcommand{\arraystretch}{1.3}
% if using array.sty, it might be a good idea to tweak the value of
% \extrarowheight as needed to properly center the text within the cells
\centering
% Some packages, such as MDW tools, offer better commands for making tables
% than the plain LaTeX2e tabular which is used here.
\begin{tabular}{|c||c|}
\hline
One & Two\\
\hline
Three & Four\\
\hline
\end{tabular}
\caption[Another example table]{Another example of a table.}
\label{tab.example_2}
\end{table}

Filing your dissertation or thesis is the last step in the process leading to the awarding of your degree. The final copy of your thesis or dissertation, which is ultimately deposited in the University Library, becomes a permanent and official record.
%
\begin{figure}[tbh]
\centering
\includegraphics{ucdavisthesis_example_figure}
% where an .eps filename suffix will be assumed under latex,
% and a .pdf suffix will be assumed for pdflatex
\caption[First sample figure]{A sample figure.}
\label{fig.sample_1}
\end{figure}

Filing your dissertation or thesis is the last step in the process leading to the awarding of your degree. The final copy of your thesis or dissertation, which is ultimately deposited in the University Library, becomes a permanent and official record.

Filing your dissertation or thesis is the last step in the process leading to the awarding of your degree. The final copy of your thesis or dissertation, which is ultimately deposited in the University Library, becomes a permanent and official record.
%
\begin{figure}[tbh]
\centering
\includegraphics[width=0.5in]{ucdavisthesis_example_figure}
% where an .eps filename suffix will be assumed under latex,
% and a .pdf suffix will be assumed for pdflatex
\caption[Second sample figure]{A reduced version of the sample figure.}
\label{fig.sample_2}
\end{figure}

Filing your dissertation or thesis is the last step in the process leading to the awarding of your degree. The final copy of your thesis or dissertation, which is ultimately deposited in the University Library, becomes a permanent and official record.

Filing your dissertation or thesis is the last step in the process leading to the awarding of your degree. The final copy of your thesis or dissertation, which is ultimately deposited in the University Library, becomes a permanent and official record.



\bibliography{ucdavisthesis_example}

% The UMI abstract uses square brackets!
\UMIabstract[The abstract that is submitted to UMI must be formatted as shown in the example here. The body of the abstract cannot exceed 350 words. It should be in typewritten form, double-spaced, and on bond paper. It is important to write an abstract that gives a clear description of the content and major divisions of the dissertation, since UMI will publish the abstract exactly as submitted. Students completing their requirements under Plan A should provide extra copies of the typed summary for use by the dissertation committee during the examination.

The abstract that is submitted to UMI must be formatted as shown in the example here. The body of the abstract cannot exceed 350 words. It should be in typewritten form, double-spaced, and on bond paper. It is important to write an abstract that gives a clear description of the content and major divisions of the dissertation, since UMI will publish the abstract exactly as submitted. Students completing their requirements under Plan A should provide extra copies of the typed summary for use by the dissertation committee during the examination.

The abstract that is submitted to UMI must be formatted as shown in the example here. The body of the abstract cannot exceed 350 words. It should be in typewritten form, double-spaced, and on bond paper. It is important to write an abstract that gives a clear description of the content and major divisions of the dissertation, since UMI will publish the abstract exactly as submitted. Students completing their requirements under Plan A should provide extra copies of the typed summary for use by the dissertation committee during the examination.

The abstract that is submitted to UMI must be formatted as shown in the example here. The body of the abstract cannot exceed 350 words. It should be in typewritten form, double-spaced, and on bond paper. It is important to write an abstract that gives a clear description of the content and major divisions of the dissertation, since UMI will publish the abstract exactly as submitted. Students completing their requirements under Plan A should provide extra copies of the typed summary for use by the dissertation committee during the examination.]

\end{document} 
